\scanpage{9}
\mspage{4}
\margin{\unsure{le donner} \uline{après} la proposition qui suit, \unsure{ou postuler}\ldots}
\uline{Proposition} Si \deleted{$R_0\ne\varnothing$}, \inserted{$\operatorname{Card}R_0\geq2$}, \deleted{\unsure{alors} \unsure{(si $R\in\mathcal M_1$)}, \unsure{$\mathcal E$} \unsure{dans}} $\mathcal E$
\deleted{\unsure{existe un} \unsure{$(\operatorname{Hom}_{\mathcal E}(e,R))$} \unsure{de l'objet} \unsure{$R_0$}}
\deleted{$\operatorname{Hom}_{\mathcal E}(\underset{\inserted{\substack{\unsure{\shortparallel|''}\\\{e\}}}}{R^0},R)\ne\varnothing$} dans \unsure{des}
\deleted{$\mathcal E$ \unsure{est}, \unsure{a des} sommes finies} \unsure{(uniquement} et $X\mapsto|X|$ y commute.

\figure{\begin{tikzpicture}\node[draw,ellipse] {$M5$};\end{tikzpicture}}
\uline{Corollaire} Sur $\forall$ ens fini il $\exists$ une structure modérée et une seule, pour laquelle les parties uniponctuelles soient modérées (str. mod. discrète).

\unsure{On se propose} le cas où on veut construire $X_1\amalg_Z X_2$, sans supposer $Z$ vide. Notons d'abord ceci :

\uline{Proposition} Soit $I$ un \deleted{type de} \inserted{type de} diagrammes fini, et $(X_i)_{i\in I}$ un diagramme de type $I$ dans $\mathcal E=\operatorname{Esp}_{\mathcal M_*}$, i.e. d'espaces modérés. Soit $X$ la limite inductive \uline{dans} $(\mathrm{Ens})$, qui est donc un \inserted{certain} quotient de $\coprod X_i$ (comme toujours). Supposons qu'il existe sur $X$ une structure modérée, telle que les appl. can. $g_i:X_i\longrightarrow X$ soient modérées. Alors $X$, muni de cette structure, est \inserted{une} $\varinjlim_I$ des $X_i$.
\mspage{8}
